% Background and Motivation: trace structure and matched physical KV costs.
\section{Background and Motivation}
\label{sec:background-motivation}

\subsection{Background}
\label{sec:motivation-background}

\paragraph{Multi-agent workflows and completion time.}
A multi-agent application completes a user task through several agents
whose results depend on one another. Each agent maintains a conversation
and alternates LLM requests with tool execution: a model response may
invoke a tool, and the returned result becomes input to a later
request~\cite{wu2023autogen}. A request is one model invocation; an
agent's task can span several requests; a workflow connects these tasks
through sequential dependencies, parallel branches and joins.
The orchestration layer determines these dependencies, while the serving
engine executes the individual requests.

Figure~\ref{fig:agent-workflow-examples} illustrates two workflows.
In code generation, a main agent delegates a CSV-export feature to API,
UI, test-authoring and documentation agents. They work from a shared
specification and return changes for integration and testing.
In deep research, a lead agent assigns technical, cost, safety and supply
investigations, then combines their evidence into a report. This follows
the lead--subagent pattern used by research
systems~\cite{anthropic2025multiagent}. Each branch can require a different
number of model and tool calls. Completing one request therefore need
not complete an agent's task, and completing one branch need not produce
the final application result.

We evaluate the application by its \emph{workflow completion time}
(WCT). For workflow $w$, released at $a_w$ with all required outputs
ready at $c_w$, $\mathrm{WCT}_w=c_w-a_w$. This elapsed time includes
model execution, tool calls, queueing and dependency waits.
In the coding example, integration waits for the required branch outputs:
a delay on the last unfinished branch postpones integration, whereas
another branch may absorb the same delay before the join. Thus the
workflow impact of a serving interruption depends on where it occurs,
as well as how long it lasts.

\begin{figure}[t]
  \centering
  \includegraphics[width=\columnwidth]{figures/agent-workflow-examples.pdf}
  \caption{Illustrative multi-agent workflows.
    Each row is one agent: $R_i$ denotes its $i$th LLM request, with
    intervening tool results carried into the next request.
    Branches have different request counts and no shared round barriers.
    The coordinator resumes after their results join.
    Positions and lengths do not encode execution time.}
  \label{fig:agent-workflow-examples}
  \Description{Two stacked left-to-right workflows. A main agent assigns
    CSV-export API, UI, test-authoring and documentation tasks in parallel,
    then integrates their outputs and runs tests. Their request counts
    are four, three, three and two, respectively. A lead researcher assigns
    technical, cost, safety and supply investigations, with four, two,
    three and two requests, before writing a report. Each branch explicitly
    unfolds its requests and intervening tools. Starts and ends are
    staggered; all edges move forward, without loops or round barriers.}
\end{figure}

\paragraph{KV state across requests and failures.}
Successive requests often extend an agent's conversation with generated
tokens, file contents or tool results. During \emph{prefill}, the model
processes input tokens and produces per-layer keys and values (KV);
\emph{decoding} extends this state as output tokens are generated.
A later request with a matching token prefix can reuse its KV under
compatible model and position semantics, computing only the new
suffix~\cite{gao2024cachedattention,zheng2024sglang}.
For example, the API agent's next request can reuse its previous context
while processing the result of a file operation. Prefix changes limit
reuse to the portion that still matches.

The conversation remains useful between requests, including while an
agent is using tools or waiting for dependencies. Retaining its KV avoids
repeating prefix computation when execution resumes. KV is derived model
state, distinct from the token history and workflow progress maintained
by the application. For conventional per-token layouts, it grows with
context length across layers and KV heads; retaining a long history can
therefore require much more memory than saving its token IDs.
Sparse attention restricts an operation to selected KV~\cite{jiang2024minference},
but omission from one operation does not establish that the state is
unnecessary for later execution.

A serving-worker failure can make this GPU-resident state unavailable,
interrupt its requests and remove serving capacity~\cite{cao2026lumen}.
A worker may span several GPUs; a compatible surviving replica has model
weights but still needs the interrupted sequence's state. We consider
fail-stop worker failures for which accepted token histories and workflow
progress survive outside the failed worker and a compatible replica is
available. Without surviving KV, that replica replays the retained tokens
to reconstruct state. A KV checkpoint preserves an already computed
prefix; recovery retrieves it and replays only the subsequent tokens
whose KV was not saved. Asynchronous copying can leave this unsaved tail
when a failure occurs~\cite{strati2024dejavu,jayakody2026ghostserve}.
The checkpoint must outlive the protected failure: same-host DRAM may
survive a GPU-worker failure but cannot protect against loss of the host.

\paragraph{The resource cost of protection.}
Keeping state recoverable consumes capacity throughout its retention
interval and bandwidth when it is copied or moved. Hierarchical caches
extend GPU HBM with host DRAM and storage, allowing more histories to be
retained~\cite{qin2025mooncake,gao2024cachedattention}. Their location
also affects when a request can execute. KV already in the recovery
executor's HBM is immediately accessible to its kernels; DRAM reuse
requires host-to-device transfer, and storage adds a read that may cross
a network. Prefetching can hide some transfer behind computation, so
exposed waiting depends on both the access path and when the state is
needed.

These choices have costs on both sides. Discarding the last reusable KV
copy frees space but leaves its state to be recomputed when needed again. Keeping
it in a slower tier avoids that reconstruction while adding retrieval
work. Reserving more fast memory for one history leaves less space for
other histories or foreground execution; background copying can also
contend with serving transfers. Consequently, protecting a request more
aggressively can impose costs on concurrent work. Two decisions must be
considered together: which histories retain a KV checkpoint, and which
checkpoints merit fast access. Their value depends on future normal reuse and on how a
recovery delay would affect the workflow's required results.

\paragraph{Existing mechanisms and protection priorities.}
Existing systems provide complementary capabilities. Hierarchical
caching reduces repeated prefix work across
turns~\cite{gao2024cachedattention,qin2025mooncake}, while KV streaming
and checkpointing reduce reconstruction after
failures~\cite{strati2024dejavu,jayakody2026ghostserve}.
Application-aware execution exposes request dependencies, and
workflow-aware caching uses upcoming agent steps to guide eviction and
prefetching~\cite{lin2024parrot,pan2025kvflow}.
Recovery coordination also considers worker load when placing checkpoints
and redirecting interrupted requests~\cite{cao2026lumen}.

Recency, next-use proximity and reconstruction cost each describe part
of this decision. None of these quantities alone determines the effect
of restoring a request sooner on workflow completion: a recent or
expensive history may belong to a branch that can wait, while another
branch gates a required result. Normal reuse also matters even when a
branch can tolerate recovery delay. The question is therefore how to
allocate finite fast-tier capacity and checkpoint storage according to
their workflow benefit, while accounting for reuse and protection costs.
We next quantify branch-duration differences and KV readiness costs,
then illustrate how workflow position changes the completion delay
caused by recovery.

\begin{figure*}[t]
  \centering
  \captionsetup[subfigure]{font=normalsize,labelfont=normalfont,
    textfont=normalfont,justification=centering,singlelinecheck=false}
  \begin{subfigure}[t]{2.48in}
    \vspace{0pt}
    \begin{minipage}[c][1.65in][c]{\linewidth}
      \centering
      \includegraphics[width=\linewidth]{figures/workflow.pdf}
    \end{minipage}
    \caption{Four-branch workflow}
    \label{fig:motivation-workflow}
  \end{subfigure}\hfill
  \begin{subfigure}[t]{1.86in}
    \vspace{0pt}
    \begin{minipage}[c][1.65in][c]{\linewidth}
      \centering
      \includegraphics[width=\linewidth]{figures/workflow-gaps.pdf}
    \end{minipage}
    \caption{Workflow-maximum gaps}
    \label{fig:motivation-branch-gap}
  \end{subfigure}\hfill
  \begin{subfigure}[t]{2.52in}
    \vspace{0pt}
    \begin{minipage}[c][1.65in][c]{\linewidth}
      \centering
      \includegraphics[width=\linewidth]{figures/recovery-cost.pdf}
    \end{minipage}
    \caption{Context and Prefill cost}
    \label{fig:motivation-recompute}
  \end{subfigure}
  \caption{Workflow structure and the cost of reusable KV state.
    (a) Four trace-derived branches with schematic Main stages A0/A5.
    Bars preserve timing; dashed tails assume an all-branch join.
    (b) ECDF of maximum branch-duration gaps per workflow: 121 workflows
    with temporally inferred fan-in; these are not failure-time slack.
    (c) Measured Qwen2.5-7B-Instruct-1M Prefill on two A100 PCIe 40\,GB GPUs
    (TP2, BF16, batch one), with sparse attention above 32K.
    Reuse starts from $(C-1)$ cached tokens and computes one;
    DRAM/shared-storage I/O is serialized with computation.
    HBM/DRAM: median and min--max of three timings;
    shared storage/full prefill: one timing per length.
    Protocol: Appendix~\ref{sec:motivation-protocol}.
    K~=~1,024; M~=~1,000,000.}

  \label{fig:motivation}
  \Description{Left: a four-branch workflow with A2 completing last at
    151.6 seconds. Middle: an ECDF of workflow-maximum branch-duration
    differences with a median of 138 seconds and a 90th percentile of
    1173 seconds. Right: four measured Prefill curves for Qwen2.5-7B-Instruct-1M
    from 8K to 1,010,000 tokens: blue circles for HBM reuse, orange squares
    for serialized DRAM reuse, purple diamonds for serialized shared-storage
    reuse, and red triangles for full prefill. Both axes are logarithmic.}
\end{figure*}

\subsection{Motivation}

\paragraph{Workflow completion requires different protection priorities.}
\label{sec:motivation-workflow}
Figure~\ref{fig:motivation}(\subref{fig:motivation-workflow}) makes the
workflow-level consequence concrete. A2 finishes last, 151.6\,s
after the fork, whereas A1 finishes at 60.4\,s. If the following Main stage
requires all four outputs, accelerating A1 alone does not advance the
join, while delaying A2 postpones downstream work. A branch can absorb
some additional delay without postponing the required output; we call
this tolerance \emph{slack}.

The public SemiAnalysis CC/Weka trace contains 393 filtered Claude Code
sessions~\cite{semianalysis2026cctraces}. We retain overlapping subagent
groups whose completed branches are followed by a Main request after the
last branch finishes; Main may also execute while the branches run.
This temporal fan-in proxy retains 286 groups in 121 workflows.
For each group we take the longest minus the shortest branch duration,
then retain each workflow's maximum gap.
Figure~\ref{fig:motivation}(\subref{fig:motivation-branch-gap}) gives a
median of 138\,s and a 90th percentile of 1,173\,s. These conditional
statistics establish substantial branch-duration heterogeneity. They do
not establish result-consumption edges, which the trace does not record,
or the slack distribution at a fault. The four-branch example illustrates
one retained group, not a representative sample of workflow maxima.

\paragraph{Long contexts make fast recovery state expensive.}
\label{sec:motivation-long-context}
Losing an agent's KV can turn a brief interruption into substantial
reconstruction work. Figure~\ref{fig:motivation}(\subref{fig:motivation-recompute})
compares four paths using Qwen2.5-7B-Instruct-1M~\cite{qwen2025qwen25oneM}
on two A100 PCIe 40\,GB GPUs, with TP2, BF16 and batch one.
All paths use the same SGLang 0.5.10rc0 model execution and
MInference vertical/slash attention~\cite{jiang2024minference} with dual chunk
attention (DCA) position handling; attention is sparse above 32K.
For total context $C=M+1$, full prefill starts from empty KV, while reuse
starts from the same $M$-token history in HBM, DRAM, or shared storage and
computes one token. Timing ends with all $C$ tokens' KV ready, before
vocabulary projection. HBM provides a resident-state reference; the
DRAM and shared-storage paths serialize restoration and computation.

At 1,010,000 tokens, full prefill takes 571.21\,s despite sparse attention.
HBM, DRAM and shared-storage reuse take 1.305, 2.663 and 10.156\,s,
respectively. A reusable history avoids most reconstruction work, but its
storage tier still changes readiness time: shared storage adds 7.49\,s
over DRAM in this run. The corresponding history occupies 57.92\,GB
across the two ranks. Keeping many such histories in DRAM can therefore
be expensive, while moving them to storage introduces additional delay.
These measurements characterize one fixed synthetic input and serialized
paths, with the timing counts stated in the caption. They establish the
cost difference in this configuration; storage-cache state and overlapping
I/O can change its magnitude. Appendix~\ref{sec:motivation-protocol}
gives the measurement and validation details.

\begin{figure}[t]
  \centering
  \includegraphics[width=\linewidth]{figures/recovery-exposure.pdf}
  \caption{Measured Prefill costs for
    Figure~\ref{fig:motivation}(\subref{fig:motivation-workflow}), using
    Figure~\ref{fig:motivation}(\subref{fig:motivation-recompute})'s configuration
    at trace input-size proxies plus one token.
    Bars/whiskers show median/min--max over three timings after warmup.
    Hatching denotes absorbed delay; solid fill is the calculated WCT
    increase under an all-branch join with other execution fixed.
    Each case is a single-branch counterfactual, not a fault-injected
    WCT measurement.}
  \label{fig:motivation-recovery-exposure}
  \Description{A1 through A4 have measured median Prefill costs of
    33.2, 32.6, 35.4 and 39.9 seconds.
    Applying each cost separately as an added branch pause gives WCT
    increases of 0.0, 32.6, 16.3 and 18.0 seconds.
    Hatching identifies the absorbed portion; whiskers show the minimum
    and maximum of three timings, not confidence intervals.}
\end{figure}

\paragraph{Failure position determines the workflow benefit of protection.}
Reconstruction cost and workflow position must be considered together.
Figure~\ref{fig:motivation-recovery-exposure} uses the four branches in
Figure~\ref{fig:motivation}(\subref{fig:motivation-workflow}). We measure
full prefill at each active request's input-size proxy, approximately
101K--114K tokens in decimal units, using the configuration above.
The proxies come from trace hash-block counts; the measured input is
synthetic, rather than the original Claude prompts.
At 29.5\,s after the fork, we insert each measured cost $R_i$ as an
additional pause on branch $i$, separately, while holding other execution
and queues fixed. Under
the assumed all-branch join, the trace completion times give offline
slack $\sigma_i$, and the resulting WCT increment is
$\max(0,R_i-\sigma_i)$.

A1's 33.2\,s cost fits within its 91.2\,s slack and adds no WCT
in this construction. A2 takes 32.6\,s but has zero slack,
exposing its entire cost. A3 and A4 expose 16.3 and
18.0\,s, respectively. Thus, reconstruction latency alone does
not determine the benefit of protection. This example combines physical
reconstruction timings with a conditional dependency calculation; it
does not measure failures in the original workflow or establish the
frequency of such outcomes. Common-worker failures and shared recovery
resources require joint analysis, developed in the next section.

\label{sec:motivation-goals}
These results motivate AegisKV's central objective: allocate finite
fast-tier capacity and backing storage according to their effect on
workflow completion. Retaining a usable copy avoids reconstruction;
keeping KV closer to execution can additionally avoid restore waits.
Neither saving has the same WCT value for every branch. Protection must also account for normal reuse
and the capacity and I/O it consumes. The design therefore separates the
value of fast-tier retention from the value of keeping any recoverable
copy, and evaluates both in the context of the workflow.
